% Part: incompleteness
% Chapter: representability-in-q
% Section: sigma1-completeness

\documentclass[../../../include/open-logic-section]{subfiles}

\begin{document}

\olfileid[id]{inc}{inp}{s1c}
\olsection{Kelengkapan \texorpdfstring{$\Sigma_1$}{Sigma-1}}

Walaupun $\Th{Q}$ dan perluasan-perluasannya yang konsisten serta dapat
diaksiomatisasi tidak lengkap, telah kita lihat bahwa $\Th{Q}$ memang
membuktikan banyak fakta dasar tentang numeral. Bahkan, hasil ini dapat
diperluas secara cukup berarti. Untuk memahami cakupan hal-hal yang dapat
dibuktikan dalam~$\Th{Q}$, kita memperkenalkan gagasan !!{formula}s
$\Delta_0$, $\Sigma_1$, dan $\Pi_1$. Secara kasar, !!{formula} $\Sigma_1$
ialah formula berbentuk $\lexists[x][!B(x)]$, dengan $!B$ dibentuk hanya
menggunakan penghubung proposisional dan kuantor terbatas. Kita akan
menunjukkan bahwa jika $!A$ adalah !!{sentence} $\Sigma_1$ yang benar dalam
$\Struct{N}$, maka $\Th{Q} \Proves !A$
(\olref{thm:sigma1-completeness}).

\begin{defn}
\ollabel{defn:bd-quant}
\emph{!!{formula} eksistensial terbatas} ialah formula berbentuk
$\lexists[x][(x < t \land !A(x))]$, dengan $t$ sebarang suku, yang secara
konvensional kita tulis sebagai $\bexists{x < t}{!A(x)}$.
%
\emph{!!{formula} universal terbatas} ialah formula berbentuk
$\lforall[x][(x < t \lif !A(x))]$, dengan $t$ sebarang suku, yang secara
konvensional kita tulis sebagai $\bforall{x < t}{!A(x)}$.
\end{defn}

\begin{defn}
\ollabel{defn:delta0-sigma1-pi1-frm}
!!{formula} $!B$ merupakan $\Delta_0$ jika dibentuk dari !!{formula}s atomik
hanya dengan menggunakan penghubung proposisional dan kuantifikasi terbatas.
%
!!{formula} $!A$ merupakan $\Sigma_1$ jika
$!A \ident \lexists[x][!B(x)]$ dengan $!B$ merupakan $\Delta_0$.
%
!!{formula} $!A$ merupakan $\Pi_1$ jika
$!A \ident \lforall[x][!B(x)]$ dengan $!B$ merupakan $\Delta_0$.
\end{defn}

\begin{lem}
\ollabel{lem:q-proves-clterm-id} Andaikan $t$ adalah suku tertutup sedemikian
sehingga $\Value{t}{N} = n$. Maka $\Th{Q} \Proves \eq[t][\num n]$.
\end{lem}

\begin{proof}
Kita membuktikannya dengan induksi pada kompleksitas~$t$. Untuk kasus dasar,
$\Value{\Obj 0}{N} = 0$, dan $\Th{Q} \Proves \eq[\Obj 0][\num 0]$ karena
$\num 0 \ident \Obj 0$.
%
Untuk kasus induktif, misalkan $t_1$ dan $t_2$ merupakan suku sedemikian
sehingga $\Value{t_1}{N} = n_1$, $\Value{t_2}{N} = n_2$,
$\Th{Q} \Proves \eq[t_1][\num n_1]$, dan
$\Th{Q} \Proves \eq[t_2][\num n_2]$.

Maka $\Value{(t_1')}{N} = n_1 + 1$, dan kita mempunyai
$\Th{Q} \Proves \eq[t_1'][{\num n_1}']$ berdasarkan aturan orde pertama
bagi identitas yang diterapkan pada hipotesis induksi dan !!{formula}
$\eq[\num{n_1}'][\num{n_1}']$, sehingga
$\Th{Q} \Proves \eq[t_1'][\num{n_1 + 1}]$ menurut definisi numeral.

Untuk penjumlahan, kita mempunyai
\[
      \Value{(t_1 + t_2)}{N}
    = \Value{t_1}{N} + \Value{t_2}{N}
    = n_1 + n_2.
\]
Menurut hipotesis induksi dan aturan identitas,
$\Th{Q} \Proves \eq[t_1 + t_2][\num{n_1} + t_2]$, lalu
$\Th{Q} \Proves \eq[t_1 + t_2][\num{n_1} + \num{n_2}]$
berdasarkan penerapan kedua aturan identitas. Menurut
\olref[inc][req][bre]{lem:q-proves-add},
$\Th{Q} \Proves \eq[\num{n_1} + \num{n_2}][\num{n_1 + n_2}]$, sehingga
$\Th{Q} \Proves \eq[t_1 + t_2][\num{n_1 + n_2}]$.

Penalaran serupa juga berlaku bagi~$\times$, dengan menggunakan
\olref[inc][req][bre]{lem:q-proves-mult}.
%
Karena ini mencakup semua suku tertutup aritmetika, kita mempunyai
$\Th{Q} \Proves \eq[t][\num n]$ bagi semua suku tertutup~$t$ sedemikian
sehingga $\Value{t}{N} = n$.
\end{proof}

\begin{prob}
Buktikan secara terperinci bagian \olref{lem:q-proves-clterm-id} yang
melibatkan~$\times$.
\end{prob}

\begin{lem}
\ollabel{lem:atomic-completeness}
Andaikan $t_1$ dan $t_2$ merupakan suku tertutup. Maka
\begin{enumerate}
\item Jika $\Value{t_1}{N} = \Value{t_2}{N}$,
    maka $\Th{Q} \Proves \eq[t_1][t_2]$.
\item Jika $\Value{t_1}{N} \neq \Value{t_2}{N}$,
    maka $\Th{Q} \Proves \eq/[t_1][t_2]$.
\item Jika $\Value{t_1}{N} < \Value{t_2}{N}$,
    maka $\Th{Q} \Proves t_1 < t_2$.
\item Jika $\Value{t_2}{N} \leq \Value{t_1}{N}$,
    maka $\Th{Q} \Proves \lnot(t_1 < t_2)$.
\end{enumerate}
\end{lem}

\begin{proof}
Diberikan suku $t_1$ dan $t_2$, tetapkan $n = \Value{t_1}{N}$ dan
$m = \Value{t_2}{N}$.

Andaikan $!A \ident t_1 = t_2$. Menurut \olref{lem:q-proves-clterm-id},
$\Th{Q} \Proves \eq[t_1][\num n]$ dan
$\Th{Q} \Proves \eq[t_2][\num m]$. Jika $n = m$, maka
$\Th{Q} \Proves \eq[\num n][\num m]$, sehingga
$\Th{Q} \Proves \eq[t_1][t_2]$ berdasarkan transitivitas identitas.
Jika $n \neq m$, maka $\Th{Q} \Proves \eq/[\num n][\num m]$, dan kembali
berdasarkan transitivitas identitas,
$\Th{Q} \Proves \eq/[t_1][t_2]$.

Sekarang misalkan $!A \ident t_1 < t_2$. Untuk kedua kasus, kita mengandalkan
aksioma~$!Q_8$, yang menyatakan bahwa
$x < y \liff \lexists[z][\eq[z' + x][y]]$ bagi semua $x,y$.

Andaikan $\Sat{N}{t_1 < t_2}$. Maka terdapat suatu $k \in \Nat$ sedemikian
sehingga $n + k + 1 = m$. Menurut \olref{lem:q-proves-clterm-id},
$\Th{Q} \Proves \eq[t_1][\num n]$ dan
$\Th{Q} \Proves \eq[t_2][\num m]$, dan menurut bagian pertama lema ini,
$\Th{Q} \Proves \eq[ {\num k}' + \num n][\num m]$. Berdasarkan
transitivitas identitas, diperoleh
$\Th{Q} \Proves \eq[ {\num k}' + t_1][t_2]$, sehingga
$\Th{Q} \Proves \lexists[z][\eq[z' + t_1][t_2]]$. Menurut arah kanan-ke-kiri
dari~$!Q_8$, $\Th{Q} \Proves t_1 < t_2$.

Sebaliknya, andaikan $\Sat/{N}{t_1 < t_2}$, yakni $m \leq n$.
%
Kita bekerja dalam~$\Th{Q}$ dan mengasumsikan bahwa $t_1 < t_2$. Menurut arah
kiri-ke-kanan dari~$!Q_8$, terdapat suatu~$z$ sedemikian sehingga
$\eq[z' + t_1][t_2]$. Karena $\Th{Q} \Proves \eq[t_1][\num n]$ dan
$\Th{Q} \Proves \eq[t_2][\num m]$, diperoleh
$\eq[z' + \num n][\num m]$.
%
Dengan induksi eksternal pada~$m$ menggunakan~$!Q_5$, diperoleh
$\eq[z' + \num{n - m}][\Obj 0]$. Jika $m = n$, maka
$\eq[z'][\Obj 0]$, yang menghasilkan kontradiksi melalui~$!Q_2$.
Jika $m < n$, maka $\eq[(z' + \num{n - m - 1})'][\Obj 0]$ menurut~$!Q_5$
sekali lagi, yang menghasilkan kontradiksi melalui~$!Q_2$.
Jadi, $\Th{Q} \Proves \lnot(t_1 < t_2)$.
\end{proof}

\begin{lem}
\ollabel{lem:bounded-quant-equiv}
Andaikan $!A$ adalah !!a{formula}, $t$ suku tertutup, dan
$k=\Value{t}{N}$. Maka
\begin{enumerate}
\item $\Th{Q} \Proves \bforall{x<t}{!A(x)}$ jika dan hanya jika
    $\Th{Q} \Proves !A(\num 0) \land \dots \land !A(\num{k-1})$.
\item $\Th{Q} \Proves \bexists{x<t}{!A(x)}$ jika dan hanya jika
    $\Th{Q} \Proves !A(\num 0) \lor \dots \lor !A(\num{k-1})$.
\end{enumerate}
\end{lem}

\begin{proof}
    Kita membuktikan kasus bagi kuantor universal terbatas. Jika
    $\Value{t}{N} = 0$, maka ruas kiri ekuivalensi dapat dibuktikan
    dalam~$\Th{Q}$ karena tidak ada $x<\num 0$ menurut
    \olref[inc][req][min]{lem:less-zero}. Demikian pula, konjungsi kosong
    dapat kita anggap sebagai $\ltrue$, yang juga dapat dibuktikan
    dalam~$\Th{Q}$.
    %
    Karena itu, andaikan $\Value{t}{N} = k+1$ untuk suatu bilangan asli~$k$.
    Menurut \olref{lem:q-proves-clterm-id}, kita dapat mengasumsikan bahwa
    kita bekerja dengan !!a{formula} berbentuk
    $\bforall{x<\num{k+1}}{!A(x)}$.

    Andaikan $\Th{Q} \Proves \bforall{x<\num{k+1}}{!A(x)}$, dan misalkan
    $n \leq k$. Karena $\Th{Q} \Proves \num n < \num{k+1}$ menurut
    \olref{lem:atomic-completeness}, berdasarkan logika diperoleh
    $\Th{Q} \Proves !A(\num n)$. Dengan menerapkan fakta ini sebanyak $k+1$
    kali, masing-masing bagi $n \leq k$, kita memperoleh
    $\Th{Q} \Proves !A(\num 0) \land \dots \land !A(\num k)$ sebagaimana
    diinginkan.

    Untuk arah lainnya, andaikan
    $\Th{Q} \Proves !A(\num 0) \land \dots \land !A(\num k)$. Dengan
    bekerja dalam $\Th{Q}$, andaikan $x < \num{k+1}$. Menurut
    \olref[inc][req][min]{lem:less-nsucc}, kita mempunyai
    $x = \num 0 \lor \dots \lor x = \num k$, sehingga berdasarkan logika
    diperoleh~$!A(x)$, dan dengan demikian klaim universal
    $\bforall{x<\num{k+1}}{!A(x)}$ mengikuti.

    Pembuktian ekuivalensi bagi !!{formula}s yang dikuantifikasi secara
    eksistensial terbatas serupa.
\end{proof}

\begin{prob}
Berikan pembuktian terperinci bagi kasus eksistensial dalam
\olref{lem:bounded-quant-equiv}.
\end{prob}

\begin{lem}
\ollabel{lem:delta0-completeness}
Jika $!A$ adalah !!{sentence} $\Delta_0$ yang benar dalam $\Struct{N}$,
maka $\Th{Q} \Proves !A$.
\end{lem}

\begin{proof}
Kita membuktikannya dengan induksi pada kompleksitas !!{formula}.
%
Kasus dasar diberikan oleh \olref{lem:atomic-completeness}, sehingga kita
beralih ke langkah induksi. Demi kesederhanaan, kita membagi kasus negasi
menjadi beberapa subkasus berdasarkan struktur !!{formula} yang dikenai
negasi.

\begin{enumerate}
\item Andaikan $(!A \land !B)$ benar dalam $\Struct{N}$, sehingga $!A$ dan
$!B$ benar dalam~$\Struct{N}$. Menurut hipotesis induksi,
$\Th{Q} \Proves !A$ dan $\Th{Q} \Proves !B$, sehingga berdasarkan logika
$\Th{Q} \Proves (!A \land !B)$.
%
\item Andaikan $\lnot (!A \land !B)$ benar dalam $\Struct{N}$, sehingga
$\lnot !A$ atau $\lnot !B$ benar dalam $\Struct{N}$. Tanpa mengurangi
keumuman, andaikan yang pertama. Menurut hipotesis induksi,
$\Th{Q} \Proves \lnot !A$, sehingga berdasarkan logika
$\Th{Q} \Proves \lnot (!A \land !B)$.
%
\item Andaikan $(!A \lor !B)$ benar dalam $\Struct{N}$, sehingga $!A$ benar
dalam $\Struct{N}$ atau $!B$ benar dalam $\Struct{N}$. Tanpa mengurangi
keumuman, andaikan yang pertama berlaku. Menurut hipotesis induksi,
$\Th{Q} \Proves !A$, sehingga berdasarkan logika
$\Th{Q} \Proves (!A \lor !B)$.
%
\item Andaikan $\lnot(!A \lor !B)$ benar dalam $\Struct{N}$, sehingga
$\lnot !A$ dan $\lnot !B$ benar dalam $\Struct{N}$. Menurut hipotesis
induksi, $\Th{Q} \Proves \lnot !A$ dan
$\Th{Q} \Proves \lnot !B$. Akibatnya, berdasarkan logika
$\Th{Q} \Proves \lnot(!A \lor !B)$.
%
\item Andaikan $\bforall{x<t}{!A(x)}$ benar dalam~$\Struct{N}$, dengan $t$
merupakan suku tertutup dan $k=\Value{t}{N}$. Menurut hipotesis induksi dan
logika, jika $!A(\num n)$ benar dalam~$\Struct{N}$ bagi semua
$n < \Value{t}{N}$, maka
$\Th{Q} \Proves !A(\num 0) \land \dots \land !A(\num{k-1})$. Menurut
\olref{lem:bounded-quant-equiv}, diperoleh
$\Th{Q} \Proves \bforall{x<t}{!A(x)}$.
%
\item Kasus bagi kuantor eksistensial terbatas, ketika kita mempunyai
!!a{sentence} berbentuk $\bexists{x < t}{!A(x)}$, serupa dengan kasus bagi
kuantor universal terbatas.
%
\item Andaikan $\lnot \bforall{x<t}{!A(x)}$ benar dalam~$\Struct{N}$, dengan
$t$ merupakan suku tertutup. !!{sentence} ini ekuivalen dengan !!{sentence}
$\bexists{x<t}{\lnot !A(x)}$, dengan ekuivalensinya dapat diturunkan dalam
$\Th{Q}$, sehingga kita dapat menerapkan penalaran bagi kuantor eksistensial
terbatas.
%
\item Demikian pula, andaikan $\lnot \bexists{x<t}{!A(x)}$ benar dalam
$\Struct{N}$, dengan $t$ merupakan suku tertutup. !!{sentence} ini ekuivalen
dalam $\Th{Q}$ dengan $\bforall{x<t}{\lnot!A(x)}$, sehingga kita dapat
menerapkan penalaran bagi kuantor universal terbatas.
%
\item Terakhir, andaikan $\lnot !A$ benar dalam $\Struct{N}$. Satu-satunya
kasus yang tersisa ialah ketika $!A$ atomik dan ketika
$\lnot !A \ident \lnot\lnot !B$ bagi suatu !!{sentence} $\Delta_0$ $!B$.
Jika $!A$ atomik, maka menurut \olref{lem:atomic-completeness},
$\Th{Q} \Proves \lnot !A$. Jika
$\lnot !A \ident \lnot\lnot !B$, maka berdasarkan logika formula itu dapat
dibuktikan ekuivalen dalam~$\Th{Q}$ dengan~$!B$, yang benar dalam
$\Struct{N}$ karena $\lnot !A$ benar dalam~$\Struct{N}$. Oleh karena itu,
menurut hipotesis induksi kita mempunyai $\Th{Q} \Proves \lnot !A$.
\end{enumerate}
\end{proof}

\begin{prob}
Berikan pembuktian terperinci bagi kasus eksistensial dalam
\olref{lem:delta0-completeness}.
\end{prob}

\begin{thm}
\ollabel{thm:sigma1-completeness}
Jika $!A$ adalah !!{sentence} $\Sigma_1$ yang benar dalam~$\Struct{N}$,
maka $\Th{Q} \Proves !A$.
\end{thm}

\begin{proof}
Jika $\lexists[x][!A(x)]$ adalah !!{sentence} $\Sigma_1$ yang benar
dalam~$\Struct{N}$, maka terdapat bilangan asli~$n$ dan penetapan
variabel~$s$ sedemikian sehingga $s(x) = n$ dan $\Sat{N}{!A(x)}[s]$.
Menurut fakta-fakta standar tentang relasi kepuasan, diperoleh
$\Sat{N}{!A(\num n)}$. Namun, $!A(\num n)$ merupakan !!{formula} $\Delta_0$,
sehingga menurut \olref{lem:delta0-completeness} kita mempunyai
$\Th{Q} \Proves !A(\num n)$, dan berdasarkan logika kita juga mempunyai
$\Th{Q} \Proves \lexists[x][!A(x)]$.
\end{proof}

\end{document}
